\section{Non-Abelian Anyons on a QCA Lattice}
\label{sec:nonabelian}

The Kitaev honeycomb model~\cite{kitaev2006} is the canonical
lattice realization of non-Abelian Ising anyons.  We show that a
QCA lattice with cells on honeycomb vertices and native
$\sigma_z \sigma_z$ couplings realizes this model with no
additional hardware requirements.

\subsection{Hamiltonian}
\begin{equation}
H = -J_x \sum_{\langle i,j\rangle_x} \sigma^x_i \sigma^x_j
    -J_y \sum_{\langle i,j\rangle_y} \sigma^y_i \sigma^y_j
    -J_z \sum_{\langle i,j\rangle_z} \sigma^z_i \sigma^z_j
\end{equation}
Each bond type is a QCA-native coupling with a local basis
rotation:
\begin{itemize}
\item $z$-bond: direct $\sigma_z \sigma_z$ (QCA-native)
\item $x$-bond: $H \cdot \sigma_z \sigma_z \cdot H$
\item $y$-bond: $R_x(\pi/2) \cdot \sigma_z \sigma_z \cdot
      R_x(-\pi/2)$
\end{itemize}

\subsection{Verification}

We verified the geometry by computing the commutator with the
plaquette operator $W_p$ (alternating Pauli string
$Y Z X Y Z X$ around each hexagon):
\begin{equation}
\max_{p} \|[H, W_p]\| = 0 \quad \text{(exact)}
\end{equation}
for $L = 2, 3$.  The ground-state energy per site in the
isotropic B-phase converges to $E_0/N = -0.7873$, matching
Kitaev's exact result in the Majorana convention.

\subsection{Topological degeneracy}

On a torus, the four flux sectors are degenerate ground states
of the gapped B-phase.  We solve each sector via Majorana
fermionization (Kitaev 2006) and measure the splitting
$\Delta E = E_{\max} - E_{\min}$ across the four sectors.

For the deeply gapped phase ($J_x = J_y = 0.1$, $J_z = 1$):
\begin{center}
\begin{tabular}{lcc}
\hline
$L$ & Split $\Delta E$ & $E_0/N$ \\
\hline
4  & $4.9 \times 10^{-4}$ & $-0.502512$ \\
6  & $5.7 \times 10^{-6}$ & $-0.502505$ \\
8  & $6.4 \times 10^{-8}$ & $-0.502505$ \\
10 & $7.0 \times 10^{-10}$ & $-0.502505$ \\
\hline
\end{tabular}
\end{center}
The splitting decays exponentially, $\Delta E \sim e^{-L/\xi}$,
confirming the four-fold topological ground-state degeneracy of
the gapped B-phase.  This is the signature of Ising anyons on
the QCA lattice.

In the gapless isotropic phase ($J_x = J_y = J_z = 1$), the
splitting is algebraic, $\Delta E \sim 1/L$, with oscillations
from Dirac points moving on and off the finite-size k-grid.
This is consistent with Kitaev's analysis and serves as a
control case: the exponential scaling appears only in the
gapped phase.

\subsection{Significance}

Combined with the toric-code hardware validation
(Sec.~\ref{sec:hardware}), this establishes QCA as a platform
for both Abelian (toric code) and non-Abelian (Ising) anyons.
Braiding of Ising anyons, combined with magic-state injection,
gives universal topological quantum computation on a substrate
no different from a classical QCA array.
